\documentclass{article}
\usepackage[UTF8]{ctex}
\usepackage{amsmath,amssymb,amsfonts}
\usepackage{enumitem}
\usepackage{geometry}
\geometry{a4paper, margin=1in}
% === 错题本排版优化（不参与题目标识） ===
\setlength{\parindent}{0pt}
\setlength{\parskip}{0.35em plus 0.10em minus 0.08em}
\setlength{\emergencystretch}{3em}
\setlist{itemsep=0.25em, topsep=0.40em, parsep=0pt}
\allowdisplaybreaks[2]
% === 错题本排版优化结束 ===
\begin{document}
\section*{上海交通大学2026年高等代数真题}
\section*{37．上海交通大学2026年高等代数试卷}

\subsection{题目1}
1．（10 分）证明：对任意相异整数 \(a_{1}, a_{2}, \cdots, a_{n}\) ，多项式

\[
f(x)=\left(x-a_{1}\right)\left(x-a_{2}\right) \cdots\left(x-a_{n}\right)-1
\]

在有理数域上不可约．

\subsection{题目2}
2．（20 分）已知线性方程组

\[
\left\{\begin{array}{l}
x_{1}+2 x_{2}+2 x_{4}=0 \\
x_{1}-2 x_{2}+(b-1) x_{3}+x_{4}=1 \\
4 x_{1}-4 x_{2}+3 x_{3}+5 x_{4}=a
\end{array}\right.
\]

解空间的维数是 2 ，求 \(a, b\) 的值并求出方程组的通解。（题目表述有误，非齐次线性方程组的解集不构成线性空间，因此不存在＂维数＂一说）

\subsection{题目3}
3．（20 分）设 \(V\) 是复数域上的有限维线性空间，\(\alpha_{1}, \alpha_{2}, \alpha_{3}, \alpha_{4}\) 是 \(V\) 的一组基，\(V\) 到 \(V\) 的线性算子 \(\varphi\)在 \(V\) 上的作用如下：

\[
\varphi\left(\alpha_{1}\right)=5 \alpha_{1}, \varphi\left(\alpha_{2}\right)=5 \alpha_{2}, \varphi\left(\alpha_{3}\right)=3 \alpha_{1}+2 \alpha_{2}+5 \alpha_{3}, \varphi\left(\alpha_{4}\right)=6 \alpha_{4} .
\]

\begin{itemize}
  \item[（1）] 求 \(\varphi\) 在基 \(\alpha_{1}, \alpha_{2}, \alpha_{3}, \alpha_{4}\) 下的矩阵 \(A\) 。
  \item[（2）] 确定 \(A\) 的若尔当标准型 \(J\) 。
  \item[（3）] 试找 \(V\) 的一组新的基，使得 \(\varphi\) 在其下的矩阵为 \(J\) 。
\end{itemize}

\subsection{题目4}
4．（20 分）假设 \(A\) 和 \(B\) 是 \(3 \times 3\) 的实矩阵，且满足条件

\[
\operatorname{det}(A)=\operatorname{det}(B)=\operatorname{det}(A+B)=\operatorname{det}(A-B)=0 .
\]

证明：对于任何实数 \(x, y \in \mathbb{R}\) ，都有 \(\operatorname{det}(x A+y B)=0\) ．

\subsection{题目5}
5．（20 分）设 \(A=\left(\begin{array}{ccc}1 & -1 & 1 \\ 1 & 1 & -1\end{array}\right)\) ．

\begin{itemize}
  \item[（1）] 计算 \(A\) 的奇异值 \(s_{1} \geq s_{2}>0\) ．
  \item[（2）] 写出 \(A\) 的奇异值分解．
\end{itemize}

\subsection{题目6}
6．（20 分）设 \(V\) 是有限维内积空间， \(\mathcal{P} \in L(V)\) 满足 \(\mathcal{P}^{2}=\mathcal{P}\) ，证明： \(\mathcal{P}\) 是某个子空间 \(U\) 上的投影算子当且仅当 \(\mathcal{P}\) 是自伴的．

\subsection{题目7}
7．（20 分）证明：实数域 \(\mathbb{R}\) 上的内积空间 \(V\) 中保持任意两个向量内积不变的变换一定是线性的，从而是正交变换。

\subsection{题目8}
8．（20 分）任取数域 \(F\) 上的一个 \(m \times n\) 的矩阵 \(A, n \times s\) 的矩阵 \(B\) 。\\
（1）证明：\(r(A B)+n \geq r(A)+r(B)\) ，此处 \(r\) 为矩阵的秩函数．\\
（2）证明：上述等式成立当且仅当 \(N(A) \subset C(B)\) ，此处 \(N(A)\) 为矩阵 \(A\) 的零空间，\(C(B)\) 为矩阵 \(B\)的列空间。
\end{document}
